\section{Programas del microcontrolador}
\label{sec:codigo}
Los programas completos se entregan en el archivo comprimido, en la carpeta
\texttt{programas\_microcontrolador/}, y están en el repositorio del grupo. La
Tabla~\ref{tab:programas} indica qué programa corresponde a cada pregunta. Los listados de este
anexo son los del programa de la pregunta~3; el módulo PID está en el Listado~\ref{lst:pid}.

\begin{table}[h]
\centering\small
\caption{Programas por pregunta (rutas dentro del comprimido).}
\label{tab:programas}
\begin{tabularx}{\linewidth}{@{}>{\raggedright\arraybackslash}p{2.6cm}>{\raggedright\arraybackslash}p{6.2cm}>{\raggedright\arraybackslash}X@{}}
\toprule
Pregunta & Archivo & Contenido \\
\midrule
3 (PID sin \texttt{delay}) & \texttt{pregunta3\_pid/main\_tarea2.c} & Lazo con \texttt{esp\_timer}, ISR del sensor y telemetría (Listado~\ref{lst:main}) \\
 & \texttt{pregunta3\_pid/pid\_controller.c/.h} & PID discreto con anti-windup (Listados~\ref{lst:pid} y \ref{lst:pidh}) \\
 & \texttt{pregunta3\_pid/estimador\_velocidad.c}, \texttt{conversion\_freno.c} & Velocidad por tiempo entre marcas y grados a pulso del servo \\
6 y 7 (planta y ensayo) & \texttt{registrador\_banco/src/aplicacion.c} & Comandos \texttt{FRENO}, \texttt{MOTOR}, \texttt{ENSAYO} y registro cada 20~ms \\
 & \texttt{registrador\_banco/src/} (resto) & Sensor, servo, motor y protocolo serie \\
 & \texttt{registrador\_banco/sdkconfig.ensayo} & Pines y calibración usados en el ensayo \\
\bottomrule
\end{tabularx}
\end{table}

\lstinputlisting[caption={Lazo de control sin \texttt{delay()} para el ESP32 (\texttt{main\_tarea2.c}).},label={lst:main}]{codigo/main_tarea2.c}

\lstinputlisting[caption={Interfaz del PID (\texttt{pid\_controller.h}).},label={lst:pidh}]{\firmware/include/pid_controller.h}
